\section{Conclusion}
In this chapter, we have explored the possibility of addressing constant delays in the state observation or in the action execution by using a non-stationary memoryless policy. 
Indeed, memoryless policies naturally induce non-stationarity of the dynamics due to the partial observation of the augmented state.
Notably, there exists an optimal memoryless policy in the space of non-stationary and Markovian policies\cite[Theorem~5.1]{derman2021acting}.

Based on this result, we design an approach able to adapt to the intra-episode non-stationarity of the process. 
Inspired by the literature on lifelong learning that considers such scenarios, we propose to learn a hyper-policy that, given time as input, samples the parameters of a policy to be queried at that time.
In order to optimise for future performance, we design an estimator for it whose foundation relies on the assumption of smooth non-stationarity.
Indeed, this estimator reuses past samples to estimate the future and exponentially discounts them as they get older.
Under smoothness conditions, we demonstrate that the bias of the estimator is bounded and that the exponential discounting can control it to some extent. 
Notably, the bias vanishes when the environment and the hyper-policy are stationary.
Two terms are added to the objective on top of the future performance estimator.
First, the estimation of the past performance for the new set of hyperparameters is added. 
This term fights catastrophic forgetting by forcing the hyper-policy to perform well on past samples. 
Second, a penalisation on the variance is added. 
It forbids excessive non-stationarity of the hyper-policy that would be symptomatic of an overfitting of the past non-stationarity of the environment and a poor generalisation in the future.
Optimising this objective by gradient optimisation yields our algorithm, POLIS.

Leveraging this algorithm, we propose a simple modification to take into account the delay. 
As in dSARSA \cite{schuitema2010control}, we shift the rewards backward in time during the estimation of the performance to realign selected policies with their outcomes.

An empirical evaluation of POLIS in different lifelong scenarios, with and without delay, has then been presented. 
The algorithm has demonstrated that it can efficiently learn the structure of the non-stationarity to adapt for the future.
POLIS also avoids extra non-stationarity where it is not needed, thanks notably to the penalisation term of the objective.
Tested on delayed tasks, POLIS demonstrated robust behaviour as the delay increases.
Yet, POLIS sometimes has a high variance in its returns on the Dam environment, compared to the stationary hyper-policy.

Future research could consider stochastic delays. 
In this case, the biggest advantage of a memoryless policy is that its input is a state in $\statespace$ and does not depend on the delay as augmented approaches would. 
Therefore, memoryless policies have an input of small and fixed dimensions $\cardinal{\statespace}$. 